\documentclass[9pt,handout]{beamer}
\usepackage[orientation=landscape,size=custom,width=16,height=9,scale=0.35,debug]{beamerposter} 
\usepackage[spanish]{babel}
\usetheme{Boadilla}
\usepackage{hyperref}

\author[]{\texorpdfstring{Leonardo Besoaín }{Author}}
\title{Estadística I}
%\setbeamercovered{transparent} 
%\logo{} 
%\institute{Facultad de Economía y Negocios\\Universidad de Chile} 
\date{} 
\subtitle{Clase sincrónica 2}
\setbeamertemplate{navigation symbols}{} 
\titlegraphic{\includegraphics[width=1.0cm]{escudo-uchile}}

\begin{document}

%%%Portada%%%
\begin{frame}
\titlepage
\end{frame}
%%%%%%%%%%%%%

%%%%%%%%%%%%%%%%%%%%%%%%%%%%%%%%%%%%%%%%%%

%%%%%%%%%%%%%%%%%%%%%%%%%%%%%%%%%%%%%%%%%%
\begin{frame}{Las apuestas de Bacanes}
\justifying
Bacanes es una página web de apuestas
que 
 asigna
números entre 0 y 1 a una gran cantidad
de eventos. 

\bigskip 

Estos números se conocen
como ``probabilidades bacanes'' y, para
un evento $A$ determinado,  se denota
por ${\rm P}_{\rm bacanes}(A)$. 

\bigskip

Para cualquier evento $A$ al que
asigna una probabilidad, 
Bacanes está dispuesta tanto a comprar
como a vender, a un precio de \$1000$\cdot{\rm P}_{\rm bacanes}(A)$,
el certificado siguiente:

\bigskip

\begin{center}
\fbox{\begin{minipage}{25em}
El dueño de este certificado puede cobrar \$1000 si ocurre $A$. 
Este certificado no tiene valor alguno si no ocurre $A$. 
% Este certificado no tiene fecha de expiración.
\end{minipage}}
\end{center}    
\end{frame}

%%%%%%%%%%%%%%%%%%%%%%%%%%%%%%%%%%%%%%%%%%
\begin{frame}{Caso 1: Bacanes no cree que $\Pr(A^c)=1-\Pr(A)$}

Para fijar ideas, suponga que Bacanes
cree que
\begin{equation}
\label{eq:Bacanes.1}
\Pr(A) + \Pr(A^c) < 1.
\end{equation}

\bigskip

Ud.\ compra un certificado de $A$ y otro de $A^c$. 

\begin{itemize}
\item Ud.\ paga: $\Pr(A)+\Pr(A^c)$ que es menos que \$1000 y 
recibe \$1000 con certeza. 
\item Su ganancia neta es: $\$1000[1-\Pr(A)-\Pr(A^c)] > 0$. 
\item Como puede realizar esta operación el número de veces que desee, puede llevar a Bacanes a la quiebra.
\end{itemize}

\bigskip

\alert{Intuición}: 
\begin{itemize}
\item La ecuación \eqref{eq:Bacanes.1} dice
que $A$ y $A^c$ están ``baratos'', luego
compro ambos. 
\end{itemize}
\end{frame}

%%%%%%%%%%%%%%%%%%%%%%%%%%%%%%%%%%%%%%%%%%

%%%%%%%%%%%%%%%%%%%%%%%%%%%%%%%%%%%%%%%%%%
\begin{frame}{Caso 1: Bacanes no cree que $\Pr(A^c)=1-\Pr(A)$}

Consideremos ahora el caso en que Bacanes cree que
\[
\Pr(A) + \Pr(A^c) > 1.
\]
En este caso, Ud.\ le vende certificados por $A$ y $A^c$ a Bacanes
y, por un argumento análogo, tiene una ganancia garantizada
de $\$1000[\Pr(A)+\Pr(A^c)-1]>0$. 

\bigskip

Como se puede vender/comprar cuantos certificados
se desee, es posible llevar a Bacanes a la quiebra. 
\end{frame}

%%%%%%%%%%%%%%%%%%%%%%%%%%%%%%%%%%%%%%%%%%

%%%%%%%%%%%%%%%%%%%%%%%%%%%%%%%%%%%%%%%%%%
\begin{frame}{Caso 2: Bacanes no cree que $\Pr(A\cap B)\leq\Pr(A)$}

Bacanes cree que 
\[
\Pr(A) < \Pr(A\cap B),
\]
lo cual sabemos es imposible (?`se acuerda de Sofía?). 

\bigskip

Ud.\ compra un certificado de $A$ y vende uno de $A\cap B$. 

\bigskip

Tres resultados posibles (lo clave es que si sucede
$A\cap B$ entonces necesariamente sucede $A$):

\begin{center}
\begin{tabular}{l cccc}
\toprule
\multicolumn{2}{c}{Eventos que suceden}  & Ingresos        por                      & Ingresos          por            & Ganancia         \\
$A$  & $B$                                       & compraventa certificados   & eventos que suceden  & neta      \\
\midrule
No    & No                                                   & $1000[\Pr(A\cap B)-\Pr(A)]$ & 0        & $1000[\Pr(A\cap B)-\Pr(A)]$\\
Si      & No                                                  & $1000[\Pr(A\cap B)-\Pr(A)]$ & 1000  & $1000[1+\Pr(A\cap B)-\Pr(A)]$\\
Si      & Si                                                    & $1000[\Pr(A\cap B)-\Pr(A)]$ & 0        & $1000[\Pr(A\cap B)-\Pr(A)]$\\
\bottomrule
\end{tabular}
\end{center}

\bigskip

Luego la combinación anterior lleva a una ganancia segura de al menos
$1000(\Pr(A\cap B)-\Pr(A))>0$.
\end{frame}

%%%%%%%%%%%%%%%%%%%%%%%%%%%%%%%%%%%%%%%%%%

%%%%%%%%%%%%%%%%%%%%%%%%%%%%%%%%%%%%%%%%%%
\begin{frame}{Caso 3: Bacanes no cree en el Axioma de Aditividad}
Sean $A$ y $B$ eventos disjuntos donde Bacanes cree que:
\[
\Pr(A\cup B) < \Pr(A) + \Pr(B). 
\]

\bigskip 

Entonces Ud.\ compra $A\cup B$ y vende
$A$ y $B$. Nuevamente, compra lo barato y vende lo caro.

\bigskip

La tabla que sigue resume por qué esta estrategia lleva a la quiebra de Bacanes. 
Lo clave es que, como $A$ y $B$ son disjuntos, si sucede uno no puede suceder el otro.

\begin{center}
% \begin{small}
\begin{tabular}{l cccc}
\resizebox{\textwidth}{!}{
\toprule
\multicolumn{2}{c}{Eventos que suceden}  & Ingresos      por                        & Ingresos        por              & Ganancia         \\
$A$  & $B$                                       & compraventa certificados   & eventos que suceden  & neta      \\
\midrule
Si    & No                                                   & $1000[\Pr(A)+\Pr(B)-\Pr(A\cup B)]$ & 0  & $1000[\Pr(A)+\Pr(B) - \Pr(A\cup B)]$\\
No   & Si                                                  &   $1000[\Pr(A)+\Pr(B)-\Pr(A\cup B)]$ & 0  & $1000[\Pr(A)+\Pr(B) - \Pr(A\cup B)$\\ 
No   & No                                                   & $1000[\Pr(A)+\Pr(B)-\Pr(A\cup B)]$ & 0 &  $1000[\Pr(A)+\Pr(B) - \Pr(A\cup B)]$\\
\bottomrule
}
\end{tabular}
% \end{small}
\end{center}
\end{frame}

%%%%%%%%%%%%%%%%%%%%%%%%%%%%%%%%%%%%%%%%%%

\begin{frame}{}

Luego la combinación anterior lleva a una ganancia segura de 
 $1000[\Pr(A)+\Pr(B) - \Pr(A\cup B)]$. 
 
 \bigskip
 
 Y como la combinación de compraventas se puede realizar
 muchas veces, usted puede llevar a Bacanes a la quiebra. 
 
 \bigskip
 
 Suponga ahora que Bacanes cree que
 \[
\Pr(A\cup B) > \Pr(A) + \Pr(B). 
\]

\bigskip

?`Cuál es la combinación de compra y venta de certificados
que lleva a Bacanes a la quiebra?
\end{frame}
%%%%%%%%%%%%%%%%%%%%%%%%%%%%%%%%%%%%%%%%

%%%%%%%%%%%%%%%%%%%%%%%%%%%%%%%%%%%%%%%
\begin{frame}{Caso 4: Bacanes no cree que $\Pr(A\cup B) = \Pr(a)+\Pr(B)-\Pr(A\cap B)$}
\begin{center}
{\Huge Se deja propuesto, solo para valientes. }
\end{center}
\end{frame}



%%%%%%%%%%%%%%%%%%%%%%%%%%%%%%%%%%%
\begin{frame}{La paradoja de las votaciones}

En su clásico libro de 1957, \emph{A Theory of Democracy}, 
con más de 30 mil citas en Google Scholar, John Downs
realiza un análisis racional de si uno debiera ir a votar o no. 

\begin{itemize}
\item Costo de ir a votar: $C>0$. Toma tiempo, gasto en locomoción, etc. 
\item Beneficio de ir a votar: Si mi voto decide la elección gano $B$,
    en caso contrario no gano nada.
\item Denotando por $p$
    la probabilidad de que mi voto decida la elección, un análisis 
    costo-beneficio de ir a votar me lleva a ir si y solo si
    \[
    pB - C > 0.
    \]
    Donde estamos usando un concepto que no hemos visto
    aun, el de ``ganancia esperada''. No obstante lo anterior, 
    es intuitivamente obvio que si uno toma la decisión
    de ir a votar de manera racional uno irá si y solo si
    la probabilidad de que el voto de uno sea decisivo
    es suficientemente alta. 
\end{itemize} 
\end{frame}

%%%%%%%%%%%%%%%%%%%%%%%%%%%%%%%%%%%%%

\begin{frame}{Aplicación: Influencia máxima}

Veremos a continuación que en la mayoría
de las elecciones la probabilidad de ser
la votante que decide la elección (votante
\alert{pivotal}) es muy baja, de modo que
una análisis costo-beneficio indica
que la gente no debiera ir a votar
en países donde el voto es voluntario. 

                    \bigskip
          
          Esto se conoce como
          la ``paradoja
          de las votaciones''
          o ``paradoja de Downs''.
          
          \bigskip
          
          A continuación calculamos la probabilidad
          de que un voto sea pivotal, aplicando 
          los resultados de conteo que vimos en el Video 4.
\end{frame}

%%%%%%%%%%%%%%%%%%%%%%%%%%%%%

\begin{frame}{Se pasa a una pregunta precisa}
Los votantes eligen entre dos candidatos A y B. 
                   La probabilidad que cada uno de ellos
                   vote por A es 0,5 y la probabilidad
                   que voten por B también es 0,5. 
                   Es decir, suponemos que no hay votos nulos o blancos y
                   todos quienes votan lo hacen por A o B.
                   
                   \bigskip
                   
                   También suponemos que 
                   cada votante decide sin que 
                   le importe cómo votan los demás. 
De manera un tanto
                   informal, podemos decir que los votantes
                   deciden independientemente una de otra. 

\bigskip
          
          Determine la probabilidad de que el voto
          de cada votante sea pivotal si votan:
          \begin{enumerate}
          \item 11 electores. 
          \item 101 electores. 
          \item 1.001 electores
          \item 10.001 electores
          \item 100.001 electores.
          \item 1.000.001 electores. 
          \end{enumerate}
\end{frame}

\begin{frame}
\frametitle{Solución}

Partimos con el caso de 11 electores,  
usted y 10 electores más, y determinamos
la probabilidad de que usted sea la votante pivotal.

\bigskip

Representamos el resultado de la votación
por una 11-tupla ordenada con As y Bs, 
donde la $i$-ésima coordenada indica
por quién votó la $i$-ésima electora. 

\bigskip

Clave: Suponemos que todas las 11-tuplas son igualmente
probables. Este supuesto se basa en que

\begin{itemize}
\item La probabilidad de votar por A o B es la misma.
\item Los votos de una electora no dependen del voto del resto.
\end{itemize}

\bigskip

Luego calcular la probabilidad que nos piden
es calcular casos favorables y casos posibles. 
\end{frame}

\begin{frame}
\frametitle{}

Consideramos las 10 votantes distintas a usted. 

\bigskip

Ud.\ será la votante pivotal si y solo si
hay 5 votos para A y 5 votos para B.

\bigskip

Luego:

\begin{itemize}
\item Espacio muestral: 10-tuplas ordenadas con As  y Bs.
\item Casos posibles: $2^{10}$. Cada votante puede votar A o B. 
\item Casos favorables: $\binom{10}{5}$. Número de formas
en que se determinan las 5 coordenadas (de la 10-tupla
ordenada) que votan por A. Una vez determinados,
también quedan determinadas las coordenadas
con los votos de B.
\end{itemize}

\bigskip

En lo que sigue dentamos por $p_n$ la probabilidad
de ser la votante pivotal cuando votan $n$ personas. 

\bigskip

El raciocinio anterior muestra que:
\[
p_{11} = \frac{\binom{10}{5}}{2^{10}} \simeq 0,246..
\]
\end{frame}

\begin{frame}
\frametitle{}

Sin $n=2k+1$ es impar, el mismo raciocinio anterior
muestra que:
\[
p_{n} = \frac{\binom{2k}{k}}{2^{2k}}.
\]

\bigskip

Lo cual lleva a:

\begin{center}
\begin{tabular}{lcccccc}
\toprule
$n$:     & 11          & 101        & 1.001       & 10.001 & 100.001 & 1.000,001 \\
\midrule
$p_n$: &  0.246   & 0.0796    &  0.0252    & 000798 & 0.00252  & 0.000798                               \\
\bottomrule
\end{tabular}
\end{center}

\bigskip

Aunque las probabilidades son pequeñas, son más grandes
de lo que la mayoría esperaría. Por ejemplo, si votan 100
personas, la probabilidad de que Ud.\ sera la votante pivotal
es de 8 por ciento. 

\bigskip

Usando la aproximación de Sterling,
 $k! \sim \sqrt{2\pi}k^{k+\frac{1}{2}}e^{-k}$
se muestra que
$p_n \simeq 1/\sqrt{n\pi}$, 
          donde $n$ es el número de votantes. 
De modo que $p_n$ cae de manera
relativamente lenta con $n$. 

\bigskip

No obstante lo anterior,  
si votan 100.000 mil personas, la probabilidad de que usted
sea el voto pivotal es de 0,25\% y si vota un millón es
de 0,08\%, es decir, una en 1250. 
\end{frame}

\begin{frame}
\frametitle{Comentarios}

        Hemos supuesto que los votantes tienen la misma
         probabilidad de votar por cualquiera de los dos
         candidatos. 
         
         \bigskip
         
         ?`En qué dirección varían las probabilidades
         calculadas (?`crecen?, ?`caen?) si
         no se cumple este supuesto?
         
         \bigskip
         
         Una forma de resolver la paradoja de las votaciones
         es notando que, aun si aceptamos que los votantes realizan
         un análisis de costo-beneficio para decidir si
         van a votar, emitir el voto pivotal no es el único
         beneficio a considerar. ?`Puede listar otros beneficios?

         \bigskip
         
         Finalmente notamos que aun cuando el modelo
         de votante racional no explica por qué la gente
         va a votar sí es un buen modelo para predecir
         cuánta gente irá a votar. 
\end{frame}

\begin{frame}
\frametitle{Los fósforos de Banach}

Un fumador lleva dos cajas de fósforos para prender sus cigarrillos
con objeto de que cuando se terminen los fósforos en una tenga
fósforos en la otra y pueda fumar mientras va a comprar
una nueva caja de fósforos. 

\bigskip

Suponga que 
la fumadora parte con dos cajas de fósforos con 50 en cada una
y que elige de cuál caja saca los fósforos al azar. 

\bigskip

?`Cuál es la probabiidad de que cuando se de
cuenta de que una caja esta vacía la otra también esté vacía?
\end{frame}

\begin{frame}
\frametitle{Solución}

\begin{itemize}
\item Espacio muestral: Todas las secuencias de 1s  y 2s, donde 1 indica que eligió
la caja 1 y 2 que eligió la caja 2. Es un espacio equiprobable. 
\item Casos posibles: $2^{100}$. 
\item Casos favorables: La secuencia tiene exactamente 50 unos
y 50 dos. En tal caso, al momento de sacar el fósforo 101, 
la fumadora se da cuenta de que las dos cajas están vacías.
El número de secuencias favorables es $\binom{100}{50}$. 
\end{itemize}

\bigskip

Luego la probabilidad que nos piden es:
\[
p = \frac{\binom{100}{50}}{2^{100}} \simeq 0,08.
\]

\bigskip

Feller, en su clásico libro de probabilidades,
plantea este problema basado en el hábito
de fumador del matemático Stefan Banach, 
por eso se le conoce como el problema
de los fósforos de Banach. 

\bigskip

Una variante acorde con los tiempos
actuales reemplaza las cajas de fósforos
por tarjetas Bip! del Transantiago.
\end{frame}




\end{document}